% Abstract: at most one page. It appears first but is written last.
\prelimchapter{Abstract}

\guide{One page that tells the reader why the problem matters, what you did
and how, what the main results were (with numbers), and what the work
contributes. Write it last, and do not cite references in it.}

Predicting where people look in videos supports applications from compression
to content design and tests our understanding of human attention. Recent video
saliency models draw their accuracy from large pretrained encoders, which they
fine-tune on small eye-tracking datasets and run once for every output frame.
This thesis asks whether a self-supervised video encoder trained by feature
prediction, \acrshort{V-JEPA}~2, is sufficient for the task without adapting it, and
whether the error of the encoder's own pretrained predictor, a learned measure
of surprise, indicates where people look.

The proposed model, SURGE-D, combines the frozen encoder with a decoder of 7.4
million parameters that builds a multi-scale pyramid from four intermediate
layers, decodes it with factorised spatio-temporal convolutions and predicts
saliency maps for sixteen frames per forward pass. Each tested design
choice was decided by a rule written in a dated log before the experiment, all
models were scored once per test set by an evaluator validated against
reference code, and differences were assessed with a paired bootstrap over
videos. The evaluation covers DHF1K, UCF Sports, Hollywood-2 with openly
licensed gaze recordings, and six audio-visual datasets.

On the DHF1K validation set, a three-seed ensemble with the \acrshort{ViT}-g encoder reached
a correlation coefficient of 0.603 and a normalised scanpath saliency of 3.49,
against 0.564 and 3.33 for SalFoM, the closest foundation-model approach, scored
by the same evaluator. It was better on all six metrics, with bootstrap intervals
that exclude zero, and remained better when SalFoM's maps were given the blur
that suits them best. A causal variant that, like SalFoM, sees no future frames
was as accurate as the offline model and beat SalFoM on the same frames by 0.029
in correlation.
With an encoder of SalFoM's size, only the ensemble was reliably better.
Offline, the model needs about one sixth of SalFoM's encoder computation per
output frame, the causal variant about three times as much. After brief
fine-tuning, the model was competitive with the best published results on UCF
Sports and exceeded the best peer-reviewed results that could be verified on
four of the six audio-visual datasets, by margins within the spread of published
evaluations; it failed on conversations, where gaze follows the speaker. Used on
its own, the predictor's error was a weaker fixation cue than simple frame
differencing and, in its raw form, anti-correlated with fixations. Encoder
pre-training and scale produced the only large gains,
whereas low-rank adaptation, a learned centre prior, deep supervision,
static-image co-training, longer clips, more training videos and audio fusion
produced none that could be measured.

The thesis contributes a frozen-encoder saliency model, a documented negative
result on predictive surprise, controlled evidence on design choices and a
transparent evaluation protocol.

\vspace{\baselineskip}
\noindent\textbf{Keywords:} video saliency prediction, eye fixations,
self-supervised learning, joint-embedding predictive architecture, frozen
encoder, predictive surprise
